\documentclass[11pt]{article}
\input{quasse/shared/preamble}
\title{QuaSSE: a guide to the model}
\begin{document}
\maketitle

This guide introduces the meaning of the model and its likelihood quantities. It is intentionally
selective. The companion \href{theory.pdf}{mathematical reference} supplies conventions and a worked
derivation; \href{explorations.pdf}{explorations} discusses possible numerical improvements.

\section{Trait-dependent diversification}
QuaSSE connects a continuously evolving trait to diversification \cite{fitzjohn2010}. A lineage
with trait value $x$ splits at rate $\lambda(x)$ and dies at rate $\mu(x)$. The rates have units of
inverse time. Trait dependence means that two lineages with different current traits can have
 different instantaneous diversification rates.

For the constant-coefficient Brownian trait model, drift $v$ gives the mean trait change per unit
time and diffusion $q$ gives its variance per unit time. Over a forward interval of length $s$,
the increment has mean $vs$ and variance $qs$. Thus $q$ is not the increment's standard deviation:
that is $\sqrt{qs}$. Both daughters inherit the parental trait at a split in this model, then evolve
independently conditional on that starting value.

\section{What is observed?}
A living lineage may be sampled at the present with probability $\rho$. Fossil observations occur
at rate $\psi$ per lineage per unit time. In the package's non-removing sampling model, observation
does not itself kill the lineage. A \emph{terminal fossil} has no observed continuation, although
its lineage may continue biologically. A \emph{sampled ancestor} has an observed continuing subtree.
Neither type should be confused with an ordinary speciation event.

Trait observations need not specify the latent trait exactly. An observation model $g(y\mid x)$
connects measured value $y$ with trait $x$. The sampling process and observation density contribute
to the likelihood along with the tree. The existing
\href{../../examples/QuaSSE_fossils_fixed_tree.md}{fixed-tree fossil example} explains the supported
BEAST representation and supplies runnable inputs.

\section{Two quantities passed through the calculation}
With $t$ denoting age before the present, $E(x,t)$ is the probability that a lineage at that age and
trait leaves no samples under the specified observation process. No samples includes both living
and fossil observations, so $E$ is not simply the probability of biological extinction.

$D(x,t)$ describes the likelihood contribution of a particular observed descendant subtree,
conditional on the ancestral lineage's trait and age. Continuous observations make it a density
contribution rather than a probability bounded by one. At the present, $E(x,0)=1-\rho$; a sampled
living tip with observed trait $y$ contributes $D(x,0)=\rho g(y\mid x)$.

Calculations proceed from observations toward ancestors, combining descendant contributions at
events. Converting these contributions into a whole-tree likelihood also requires specifying the
starting trait treatment and what events are conditioned on. For example, starting at a root split
is different from starting with a single lineage before that split. This guide does not imply that
all such alternatives are implemented; the fossil example documents its root-conditioned analysis.

\section{Model versus numerical approximation}
The continuous-trait model does not itself specify a finite trait grid or a timestep. The package
approximates its likelihood using discretized trait space and numerical propagation. Domain width,
spatial spacing, timestep, and boundary treatment can each affect the answer. Refining one does not
eliminate errors due to the others.

The log-space exploration asks whether a different representation can support higher spatial order
without negative represented likelihoods or no-sample probabilities. It is an investigation, not a
new production solver or a general accuracy guarantee. Existing
\href{../../validation/FossilSampling.md}{fossil validation} and
\href{../../validation/QuaSSEReference.md}{reference comparisons} contain numerical evidence.

\bibliographystyle{plain}
\bibliography{references}
\end{document}
